\section{AMC 8 1988}
        \begin{enumerate}
        	\item (\textit{AMC 8 1988}) The diagram shows part of a scale of a measuring device.  The arrow indicates an approximate reading of $m$. What is the value of $m$?


 
                    \begin{figure}[H]
                        \centering
                        \includegraphics[width=0.5\textwidth]{Problem/AMC_8_AJHSME_1988_1_Problem_Figure_1.png}
                    \end{figure}
                    


 $\text{(A)}\ 10.05 \qquad \text{(B)}\ 10.15 \qquad \text{(C)}\ 10.25 \qquad \text{(D)}\ 10.3 \qquad \text{(E)}\ 10.6$


 
 \textbf{Answer: }D



\textbf{Solution}

Since the scale is divided into $4$ equal sections, the marks to the left and right the arrow must be $10.25$ and $10.5$, respectively.  The only answer choice strictly between $10.25$ and $10.5$ is $\boxed{\textbf{(D)}\ 10.3}$.

 


	\item (\textit{AMC 8 1988}) The product $8\times .25\times 2\times .125 =$


 $\text{(A)}\ \frac18 \qquad \text{(B)}\ \frac14 \qquad \text{(C)}\ \frac12 \qquad \text{(D)}\ 1 \qquad \text{(E)}\ 2$


 
 \textbf{Answer: }C



\textbf{Solution}

Converting decimals to fractions yields $8 \times .25 \times 2 \times .125 = \frac{8}{1} \times \frac{1}{4} \times \frac{2}{1} \times \frac{1}{8} = \frac{16}{32} = \boxed{\textbf{(C)}\ \frac{1}{2}}$.

 


	\item (\textit{AMC 8 1988}) Evaluate the following expression:


 $\frac{1}{10}+\frac{2}{20}+\frac{3}{30} =$


 $\text{(A)}\ .1 \qquad \text{(B)}\ .123 \qquad \text{(C)}\ .2 \qquad \text{(D)}\ .3 \qquad \text{(E)}\ .6$


 
 \textbf{Answer: }D



\textbf{Solution}

Each fraction simplifies to $\frac{1}{10}$, so this sum is $\frac{1}{10}+\frac{1}{10}+\frac{1}{10} = \frac{3}{10} = \boxed{\textbf{(D)}\ .3}$.

 


	\item (\textit{AMC 8 1988}) The figure consists of alternating light and dark squares.  The number of dark squares exceeds the number of light squares by


 $\text{(A)}\ 7 \qquad \text{(B)}\ 8 \qquad \text{(C)}\ 9 \qquad \text{(D)}\ 10 \qquad \text{(E)}\ 11$


 
                    \begin{figure}[H]
                        \centering
                        \includegraphics[width=0.5\textwidth]{Problem/AMC_8_AJHSME_1988_4_Problem_Figure_1.png}
                    \end{figure}
                    


 
 \textbf{Answer: }B



\textbf{Solution 1}

Notice that in each row, the number of black squares is $1$ more than the number of white squares. Since this figure has $8$ rows, it has $\boxed{\textbf{(B)}\ 8}$ more black squares than white squares.

 



\textbf{Solution 2}

Notice that after the first row, each row has $1$ more white square and $1$ more black square than the previous row.  So, this figure has $1 + 2 + \cdots + 7 + 8$ black squares and $1 + 2 + \cdots + 7$ white squares. Therefore, there are $(1 + 2 + \cdots + 7 + 8) - (1 + 2 + \cdots + 7) = \boxed{\textbf{(B)}\ 8}$ more black squares than white squares.

 


	\item (\textit{AMC 8 1988}) If $\angle \text{CBD}$ is a right angle, then this protractor indicates that the measure of $\angle \text{ABC}$ is approximately


 
                    \begin{figure}[H]
                        \centering
                        \includegraphics[width=0.5\textwidth]{Problem/AMC_8_AJHSME_1988_5_Problem_Figure_1.png}
                    \end{figure}
                    


 $\text{(A)}\ 20^\circ \qquad \text{(B)}\ 40^\circ \qquad \text{(C)}\ 50^\circ \qquad \text{(D)}\ 70^\circ \qquad \text{(E)}\ 120^\circ$


 
 \textbf{Answer: }C



\textbf{Solution}

$\angle OBA = 20^{\circ}$, $\angle CBD = 90^{\circ}$, and $\angle OBD = 160^{\circ}$.  So $\angle ABC = 160^{\circ} - 20^{\circ} - 90^{\circ} = \boxed{\textbf{(C)}\ 50^{\circ}}$.

 


	\item (\textit{AMC 8 1988}) $\frac{(.2)^3}{(.02)^2} =$


 $\text{(A)}\ .2 \qquad \text{(B)}\ 2 \qquad \text{(C)}\ 10 \qquad \text{(D)}\ 15 \qquad \text{(E)}\ 20$


 
 \textbf{Answer: }E



\textbf{Solution 1}

Converting the decimals to fractions yields $\frac{(.2)^3}{(.02)^2} = \left( \frac{1}{5}\right)^3 \div \left(\frac{1}{50}\right)^2 =  \frac{1}{125} \div \frac{1}{2500} = \frac{1}{125} \cdot \frac{2500}{1} = \frac{2500}{125}=\boxed{\textbf{(E)}\ 20}$.

 



\textbf{Solution 2}

Expanding yields $\frac{(.2)^3}{(.02)^2} = \frac{.2 \times .2 \times .2}{.02 \times .02} = \frac{.2}{.02}\times\frac{.2}{.02}\times{.2} = 10 \times 10 \times .2 = \boxed{\textbf{(E)}\ 20}$.

 



\textbf{Solution 3}

Using scientific notation, $\frac{(0.2)^3}{(0.02)^2} = \frac{(2 \times 10^{-1})^3}{(2 \times 10^{-2})^2} = \frac{8 \times 10^{-3}}{4 \times 10^{-4}} = 2 \times 10 = \boxed{\textbf{(E)}\ 20}$.

 


	\item (\textit{AMC 8 1988}) $2.46\times 8.163\times (5.17+4.829)$ is closest to


 $\text{(A)}\ 100 \qquad \text{(B)}\ 200 \qquad \text{(C)}\ 300 \qquad \text{(D)}\ 400 \qquad \text{(E)}\ 500$


 
 \textbf{Answer: }B



\textbf{Solution 1}

$2.46 \times 8.163 \times (5.17 + 4.829) = 2.46 \times 8.163 \times 9.999 \approx 2.5 \times 8 \times 10 = \boxed{\textbf{(B)}\ 200}$.

 - GaloisTorrent, edited by j314andrews,

 


	\item (\textit{AMC 8 1988}) Betty used a calculator to find the product $0.075 \times 2.56$.  She forgot to enter the decimal points.  The calculator showed $19200$.  If Betty had entered the decimal points correctly, the answer would have been


 $\text{(A)}\ .0192 \qquad \text{(B)}\ .192 \qquad \text{(C)}\ 1.92 \qquad \text{(D)}\ 19.2 \qquad \text{(E)}\ 192$


 
 \textbf{Answer: }B



\textbf{Solution 1}

Betty calculated $75 \times 256 = 19200$, when she should have calculated $\frac{75}{1000}\times\frac{256}{100} = \frac{19200}{100000} = \frac{192}{1000} = \boxed{\textbf{(B)}\ .192}$.

 



\textbf{Solution 2}

Since $0.075$ has $3$ decimal places and $2.56$ has $2$ decimal places, the decimal point in the product should be moved $3+2=5$ places to the left.  Therefore, the correct product is $.19200 = \boxed{\textbf{(B)}\ .192}$.  

 


	\item (\textit{AMC 8 1988}) An isosceles triangle is a triangle with two sides of equal length.  How many of the five triangles on the square grid below are isosceles?


 
                    \begin{figure}[H]
                        \centering
                        \includegraphics[width=0.5\textwidth]{Problem/AMC_8_AJHSME_1988_9_Problem_Figure_1.png}
                    \end{figure}
                    


 $\text{(A)}\ 1 \qquad \text{(B)}\ 2 \qquad \text{(C)}\ 3 \qquad \text{(D)}\ 4 \qquad \text{(E)}\ 5$


 
 \textbf{Answer: }D



\textbf{Solution}

Going from top left to bottom right, the first triangle has two sides of length $\sqrt{5}$ by the Pythagorean theorem, so it is isosceles.

 The second triangle has two sides of length $2$, so it is isoceles.

 The third triangle has sides of length $2$, $\sqrt{5}$, and $\sqrt{13}$, so it is not isosceles.

 The fourth triangle has two sides of length $\sqrt{10}$ so it is isosceles.

 The fifth triangle has two sides of length $\sqrt{5}$, so it is isosceles.

 So, $\boxed{\textbf{(D)}\ 4}$ of these triangles are isosceles.

 


	\item (\textit{AMC 8 1988}) Chris's birthday is on a Thursday this year.  What day of the week will it be $60$ days after her birthday?


 $\text{(A)}\ \text{Monday} \qquad \text{(B)}\ \text{Wednesday} \qquad \text{(C)}\ \text{Thursday} \qquad \text{(D)}\ \text{Friday} \qquad \text{(E)}\ \text{Saturday}$


 
 \textbf{Answer: }A



\textbf{Solution 1}

Since $60$ days is $8$ weeks and $4$ days, the day $60$ days after Chris' birthday will be $4$ days after a Thursday, which is $\boxed{\textbf{(A)}\ \text{Monday}}$.

 



\textbf{Solution 2 (Modular Arithmetic)}

Since $60\equiv4\pmod7$, the day $60$ days after Chris' birthday will be $4$ days after a Thursday, which is $\boxed{\textbf{(A)}\ \text{Monday}}$.

 


	\item (\textit{AMC 8 1988}) $\sqrt{164}$ is 


 $\text{(A)}\ 42 \qquad \text{(B)}\ \text{less than }10 \qquad \text{(C)}\ \text{between }10\text{ and }11 \qquad \text{(D)}\ \text{between }11\text{ and }12 \qquad \text{(E)}\ \text{between }12\text{ and }13$


 
 \textbf{Answer: }E



\textbf{Solution}

Since $12^2 = 144 < 164 < 169 = 13^2$, it follows that $12 < \sqrt{164} < 13$.  That is, $\sqrt{164}$ is $\boxed{\textbf{(E)}\ \text{between }12\text{ and }13}$.

 


	\item (\textit{AMC 8 1988}) Suppose the estimated $20$ billion dollar cost to send a person to the planet Mars is shared equally by the $250$ million people in the U.S. Then each person's share is 


 $\text{(A)}\ 20\text{ dollars} \qquad \text{(B)}\ 50\text{ dollars} \qquad \text{(C)}\ 80\text{ dollars} \qquad \text{(D)}\ 100\text{ dollars} \qquad \text{(E)}\ 125\text{ dollars}$


 
 \textbf{Answer: }C



\textbf{Solution 1}

The cost per person is $\frac{\$20\text{ billion}}{250\text{ million}} = \frac{\$20000\text{ million}}{250\text{ million}} = \boxed{\textbf{(C)}\ \$80}$.

 



\textbf{Solution 2}

The cost per person is $\frac{\$20\text{ billion}}{250\text{ million}} = \frac{\$20 \times 10^9}{250 \times 10^6} = \frac{\$20 \times 10\textasciicircum{}3}{250} = \frac{\$20000}{250} = \boxed{\textbf{(C)}\ \$80}$.

 


 


	\item (\textit{AMC 8 1988}) If rose bushes are spaced about $1$ foot apart, approximately how many bushes are needed to surround a circular patio whose radius is $12$ feet?


 $\text{(A)}\ 12 \qquad \text{(B)}\ 38 \qquad \text{(C)}\ 48 \qquad \text{(D)}\ 75 \qquad \text{(E)}\ 450$


 
 \textbf{Answer: }D



\textbf{Solution}

The circumference of the patio is $2\pi \cdot 12 = 24\pi \approx 24 \cdot 3.14 = 75.36$.  Since the bushes are spaced about $1$ foot apart, about $\boxed{\textbf{(D)}\ 75}$ are needed.

 


	\item (\textit{AMC 8 1988}) $\diamondsuit$ and $\Delta$ are whole numbers and $\diamondsuit \times \Delta =36$.  The largest possible value of $\diamondsuit + \Delta$ is 


 $\text{(A)}\ 12 \qquad \text{(B)}\ 13 \qquad \text{(C)}\ 15 \qquad \text{(D)}\ 20\ \qquad \text{(E)}\ 37$


 
 \textbf{Answer: }E



\textbf{Solution 1}

Since $\diamondsuit$ and $\Delta$ can be swapped to get the same sum, only cases where $\diamondsuit \geq \Delta$ need to be considered.  These cases are shown in the table:

 
 \[\begin{array}{|c|c|c|} \multicolumn{3}{}{} \\ \hline \diamondsuit & \Delta & \diamondsuit + \Delta \\ \hline 36 & 1 & 37 \\ \hline 18 & 2 & 20 \\ \hline 12 & 3 & 15 \\ \hline 9 & 4 & 13 \\ \hline 6 & 6 & 12 \\ \hline \end{array}\]

 So $\boxed{\textbf{(E)}\ 37}$ is the largest possible sum.

 



\textbf{Solution 2}

We notice that the function $f(x) = x + \frac{36}x$ is convex with a minimum at $x = 6$, so we seek to pick a value of $x$ far from 6. To ensure that the answers are positive integers, we let $x = 36$, from which we find that $f(x) = 37$, the maximum.

 -GaloisTorrent

 



\textbf{Solution 3 (Answer Choices)}

Notice that $1 \times 36 = 36$ and $1+36=37$.  Since $\boxed{\textbf{(E)}\ 37}$ is the highest answer choice it must be correct.

 


	\item (\textit{AMC 8 1988}) The reciprocal of $\left( \frac{1}{2}+\frac{1}{3}\right)$ is 


 $\text{(A)}\ \frac{1}{6} \qquad \text{(B)}\ \frac{2}{5} \qquad \text{(C)}\ \frac{6}{5} \qquad \text{(D)}\ \frac{5}{2} \qquad \text{(E)}\ 5$


 
 \textbf{Answer: }C



\textbf{Solution}

$\frac{1}{2} + \frac{1}{3} = \frac{5}{6}$, and the reciprocal of $\frac{5}{6}$ is $\boxed{\textbf{(C)}\ \frac{6}{5}}$.

 


	\item (\textit{AMC 8 1988}) Placing no more than one $\text{X}$ in each small square, what is the greatest number of $\text{X}$'s that can be put on the grid shown without getting three $\text{X}$'s in a row vertically, horizontally, or diagonally?


 $\text{(A)}\ 2 \qquad \text{(B)}\ 3 \qquad \text{(C)}\ 4 \qquad \text{(D)}\ 5 \qquad \text{(E)}\ 6$


 
                    \begin{figure}[H]
                        \centering
                        
\includegraphics[width=0.5\textwidth]{Problem/AMC_8_AJHSME_1988_16_Problem_Figure_1.png}
                    \end{figure}
                    


 
 \textbf{Answer: }E



\textbf{Solution}

It is impossible to place $7$ or more X's, as there would be $2$ or fewer empty squares and thus not every row could have an empty square.

 It is possible to place $\boxed{\textbf{(E)}\ 6}$ X's as follows:

 
                    \begin{figure}[H]
                        \centering
                        \includegraphics[width=0.5\textwidth]{Solution/AMC_8_AJHSME_1988_16_Solution_Figure_1.png}
                    \end{figure}
                    

 


	\item (\textit{AMC 8 1988}) The shaded region formed by the two intersecting perpendicular rectangles, in square units, is


 
                    \begin{figure}[H]
                        \centering
                        \includegraphics[width=0.5\textwidth]{Problem/AMC_8_AJHSME_1988_17_Problem_Figure_1.png}
                    \end{figure}
                    


 $\text{(A)}\ 23 \qquad \text{(B)}\ 38 \qquad \text{(C)}\ 44 \qquad \text{(D)}\ 46 \qquad \text{(E)}\ \text{unable to be determined from the information given}$


 
 \textbf{Answer: }B



\textbf{Solution}

The area of the $2 \times 10$ rectangle is $20$, while the area of the $8 \times 3$ rectangle is $24$.  If these two areas were added together, however, the area of the intersection, which is a $2 \times 3$ rectangle, would be counted twice.  Since the intersection has area $6$, the area of this figure must be $20 + 24 - 6 = \boxed{\textbf{(B)}\ 38}$.

 


	\item (\textit{AMC 8 1988}) The average weight of $6$ boys is $150$ pounds and the average weight of $4$ girls is $120$ pounds.  The average weight of the $10$ children is 


 $\text{(A)}\ 135\text{ pounds} \qquad \text{(B)}\ 137\text{ pounds} \qquad \text{(C)}\ 138\text{ pounds} \qquad \text{(D)}\ 140\text{ pounds} \qquad \text{(E)}\ 141\text{ pounds}$


 
 \textbf{Answer: }C



\textbf{Solution}

The total weight of the boys is $6 \times 150 = 900$ pounds, and the total weight of the girls is $4 \times 120 = 480$ pounds.  Therefore, the total weight of all $10$ children is $900 + 480 = 1380$ pounds, and their average weight is $1380 \div 10 = \boxed{\textbf{(C)}\ 138}$ pounds.

 


	\item (\textit{AMC 8 1988}) What is the $100\text{th}$ number in the arithmetic sequence: $1,5,9,13,17,21,25,...$?


 $\text{(A)}\ 397 \qquad \text{(B)}\ 399 \qquad \text{(C)}\ 401 \qquad \text{(D)}\ 403 \qquad \text{(E)}\ 405$


 
 \textbf{Answer: }A



\textbf{Solution 1}

The common difference of this arithmetic sequence is $5 - 1 = 4$.  From the first term, it takes $99$ steps to reach the $100\text{th}$ term, so the $100\text{th}$ term is $1 + 4\cdot99 = \boxed{\textbf{(A)}\ 397}$

 



\textbf{Solution 2}

As in solution 1, the common difference is $4$.  Recall that the $n\text{th}$ term of an arithmetic sequence with first term $a$ and common difference $d$ is $a + d(n-1) = 1 + 4(100-1) = \boxed{\textbf{(A)}\ 397}$.

 


	\item (\textit{AMC 8 1988}) The glass gauge on a cylindrical coffee maker shows that there are $45$ cups left when the coffee maker is $36\%$ full.  How many cups of coffee does it hold when it is full?


 $\text{(A)}\ 80 \qquad \text{(B)}\ 100 \qquad \text{(C)}\ 125 \qquad \text{(D)}\ 130 \qquad \text{(E)}\ 262$


 
                    \begin{figure}[H]
                        \centering
                        \includegraphics[width=0.5\textwidth]{Problem/AMC_8_AJHSME_1988_20_Problem_Figure_1.png}
                    \end{figure}
                    


 
 \textbf{Answer: }C



\textbf{Solution}

Let $x$ be the number of cups of coffee that the coffee maker can hold when it is full.  Then $.36x=45$, so $x= \boxed{\textbf{(C)}\ 125}$.

 


	\item (\textit{AMC 8 1988}) A fifth number, $n$, is added to the set $\{ 3,6,9,10 \}$ to make the mean of the set of five numbers equal to its median.  The number of possible values of $n$ is


 $\text{(A)}\ 1 \qquad \text{(B)}\ 2 \qquad \text{(C)}\ 3 \qquad \text{(D)}\ 4 \qquad \text{(E)}\ \text{more than }4$


 
 \textbf{Answer: }C



\textbf{Solution}

The mean of the set $\{3, 6, 9, 10, n\}$ is $\frac{3+6+9+10+n}{5} = \frac{28+n}{5}$.

 If $n < 6$, the median is $6$.  Solving $\frac{28+n}{5} = 6$ yields $n = 2$, which is valid since $2 < 6$.

 If $6 \leq n < 9$, the median is $n$.  Solving $\frac{28+n}{5} = n$ yields $n = 7$, which is valid since $6 \leq 7 < 9$.

 If $n \geq 9$, the median is $9$.  Solving $\frac{28+n}{5} = 9$ yields $n = 17$, which is valid since $17 \geq 9$.

 Therefore, there are $\boxed{\textbf{(C)}\ 3}$ possible values for $n$.

 


	\item (\textit{AMC 8 1988}) Tom's Hat Shoppe increased all original prices by $25\%$.  Now the shoppe is having a sale where all prices are $20\%$ off these increased prices.  Which statement best describes the sale price of an item?


 $\text{(A)}\ \text{The sale price is }5\% \text{ higher than the original price.}$


 $\text{(B)}\ \text{The sale price is higher than the original price, but by less than }5\% .$


 $\text{(C)}\ \text{The sale price is higher than the original price, but by more than }5\% .$


 $\text{(D)}\ \text{The sale price is lower than the original price.}$


 $\text{(E)}\ \text{The sale price is the same as the original price.}$


 
 \textbf{Answer: }E



\textbf{Solution 1}

Let $x$ be the original price of the item.  The shoppe originally increased this to $1.25x$.  The sale brings it down to $(0.8)(1.25x)=x$, so $\boxed{\textbf{(E)}\ \text{The sale price is the same as the original price.}}$

 



\textbf{Solution 2}

Suppose the original price of an item was $\$100$.  The shoppe originally increased the price to $1.25 \cdot \$100 = \$125$, while the sale reduced it to $0.8 \cdot \$125 = \$100$, so $\boxed{\textbf{(E)}\ \text{The sale price is the same as the original price.}}$

 


	\item (\textit{AMC 8 1988}) Maria buys computer disks at a price of $4$ for $<dollar>5$ and sells them at a price of $3$ for $<dollar>$$5$.  How many computer disks must she sell in order to make a profit of $<dollar>100$?


 $\text{(A)}\ 100 \qquad \text{(B)}\ 120 \qquad \text{(C)}\ 200 \qquad \text{(D)}\ 240 \qquad \text{(E)}\ 1200$


 
 \textbf{Answer: }D



\textbf{Solution 1}

Maria buys at $\frac{5}{4}$ dollars per disk and sells at $\frac{5}{3}$ dollars per disk, so her profit is $\frac{5}{3}-\frac{5}{4}=\frac{5}{12}$ dollars per disk.  Thus to make a profit of $\$100$ she must sell $100 \div \frac{5}{12} = \boxed{\textbf{(C)}\  240}$ disks.  This is valid since $240$ is divisible by both $3$ and $4$, and therefore Maria will not have any disks left over.

 



\textbf{Solution 2}

If Maria buys $12$ disks, she buys $3$ sets of $4$ so she pays $\$15$.  If Maria sells $12$ disks, she sells $4$ sets of $3$ so she gets $\$20$.  Therefore, for every $12$ disks Maria sells, she profits $\$5$.  

 To profit $\$100$, she must sell $\frac{100}{5}=20$ batches of $12$ disks, for a total of $\boxed{\textbf{(C)}\  240}$ disks.  

 


	\item (\textit{AMC 8 1988}) 
                    \begin{figure}[H]
                        \centering
                        \includegraphics[width=0.5\textwidth]{Problem/AMC_8_AJHSME_1988_24_Problem_Figure_1.png}
                    \end{figure}
                    


 The square in the first diagram "rolls" clockwise around the fixed regular hexagon until it reaches the bottom.  In which position will the solid triangle be in diagram $4$?


 
                    \begin{figure}[H]
                        \centering
                        
\includegraphics[width=0.5\textwidth]{Problem/AMC_8_AJHSME_1988_24_Problem_Figure_2.png}
                    \end{figure}
                    


 
 \textbf{Answer: }A



\textbf{Solution 1}

In this diagram, the side containing the base of the triangle is shown in red.  After $4$ steps, it will be on the left side of the square, so $\boxed{\textbf{(A)}}$ is correct.

 
                    \begin{figure}[H]
                        \centering
                        \includegraphics[width=0.5\textwidth]{Solution/AMC_8_AJHSME_1988_24_Solution_Figure_1.png}
                    \end{figure}
                    

 



\textbf{Solution 2}

In each step, the square rotates by an angle of $150^{\circ}$ clockwise.  Therefore, in $3$ steps, the square rotates $450^{\circ}$ clockwise, which is $90^{\circ}$ more than a full rotation.  Since the triangle originally pointed upward, after rotating $90^{\circ}$ clockwise it will point rightward.  Therefore, $\boxed{\textbf{(A)}}$ is correct.

 


	\item (\textit{AMC 8 1988}) A \textbf{palindrome} is a whole number that reads the same forwards and backwards.  If one neglects the colon, certain times displayed on a digital watch are palindromes.  Three examples are: $\boxed{1:01}$, $\boxed{4:44}$, and $\boxed{12:21}$.  How many times during a $12$-hour period will be palindromes?


 $\text{(A)}\ 57 \qquad \text{(B)}\ 60 \qquad \text{(C)}\ 63 \qquad \text{(D)}\ 90 \qquad \text{(E)}\ 93$


 
 \textbf{Answer: }A



\textbf{Solution}

Consider two cases:

 Case 1: the hour is $9$ or less.  In this case, the time has $3$ digits, so it is in the form $A:BA$.  Since $A$ can be any digit from $1$ to $9$ while $B$ can be any digit from $0$ to $5$, there are $9\cdot 6 = 54$ three-digit palindromes.

 Case 2: the hour is $10$ or more.  In this case, the time has $4$ digits, so it is in the form $AB:BA$.  Since $AB$ can be $10$, $11$, or $12$, there are $3$ four-digit palindromes.

 Therefore, there are $54 + 3 = \boxed{\textbf{(A)}\ 57}$ total palindromes.

 


\end{enumerate}